% ch-policy.tex -- what a version number promises, and what may change between versions
\chapter{バージョンと更新}
\label{ch:policy}

Yuzenの現在のバージョンは1.0である。本章では、バージョン番号の意味と、リリースが変わっても
何を当てにしてよいかを述べる。ソルバーにとって最も重要な、変更によって数値が変わったときの扱いもここで定める。
本章は仕組みの説明ではなく方針の表明である。コードで強制している事項はそう明記し、
強制していない事項もそう明記する。

\textbf{バージョン1.0は最初の公開リリースである}（\cref{ch:revisions}）。

\section{バージョン番号}
\label{sec:policy:scheme}

バージョンは\code{MAJOR.MINOR.PATCH}の形式である。\file{yuzen/src/Version.h}の三つのマクロで定義し、
\code{C\_Yuzen::Version()}で文字列に整形する。各レベルは利用者に対する次の約束を表す。

\begin{description}

  \item[PATCH（パッチ）] バグ修正と文書の更新。\textbf{既存の入力ファイルは変更なしで動く。}
    数値結果が変わるのは、以前の挙動が誤っていた箇所\emph{だけ}である。そうした変更はすべて、差の大きさと
    測定に用いたケースを添えて\cref{ch:revisions}に記載する。何も告知せずに数値を変えるパッチリリースは
    リリースの欠陥であり、許されない。

  \item[MINOR（マイナー）] 機能の追加。新しい入力キー、既存キーで新たに使える値、新しい物理項、新しい
    出力列などである。\textbf{既存の入力ファイルは引き続き動く}。新しい値を受け付けるようになったキーや、
    以前はなかった省略可能なキーがあっても、何も壊れない。書かれていないキーは既定値をとるからである。
    既存のケースの結果も変わらない。\cref{sec:repro:order}に示すとおり、新しい項は登録リストの
    末尾に追加され、既定では無効である。そのため、新しい項を使わないケースの結果はビット単位で一致する。

  \item[MAJOR（メジャー）] 既存の入力ファイルや結果を無効にしうる変更。キーの意味や既定値の
    変更、離散化の置き換えなどである。こうした変更はメジャーリリースでのみ行う。メジャーリリースに
    移行する際は、重要な計算をすべてやり直す前提で臨むべきである。

\end{description}

コードが実際に実装している規則は次のとおりである。
\code{C\_Yuzen::VersionSeries()}は\code{MAJOR.MINOR}だけを返し、
\file{yuzen/src/IO/Input.cpp}の\code{Input::ReadRoot()}は\key{SolverName}を受理可能な系列の
リストと照合する。したがって、入力ファイルの\key{SolverName}には\code{MAJOR.MINOR}を書き、パッチ
レベルは書かない。バージョン1.0.0ではリストの項目は\val{"Yuzen-1.0"}だけであり、それ以外の値を書くと
\code{Invalid value. SolverName must be one of [...]}を出して実行を終了する。このメッセージにはリストの内容が
表示される。

\textbf{そのため、パッチリリースで入力ファイルの書き換えが必要になることはない}。これがこの規則の要点である。

\begin{note}[マイナーリリースで古い入力ファイルが動き続ける仕組み]
  照合の相手は、\code{Input::ReadRoot()}が組み立てる受理系列の\emph{リスト}である。1.0では
  リストの項目は\val{"Yuzen-1.0"}だけである。キーの解釈を変えないマイナーリリースでは、新しい系列を
  リストに\textbf{追加}し、古い項目は残す。そのため、1.0向けに書いた入力ファイルは1.1でも
  そのまま動く。リストから項目を削るのは、キーの名前、既定値、意味が変わるなど、キーの解釈が本当に変わった
  場合に限る。系列を外すのは意図的な措置であり、古い入力ファイルが黙って別の答えを出すのではなく、
  はっきり失敗するようにするためである。項目を削った場合は必ず\cref{ch:revisions}に
  記録する。
\end{note}

バージョンだけではビルドを特定できない。\file{JobInfo.txt}には、バージョン、ビルド日、コンパイラが、
入力読み込み、格子、ユーティリティの各層のバージョンとともに記録される（\cref{sec:output:jobinfo}）。
そのため結果ディレクトリを見れば、パッチレベルが違う場合はもちろん、同じパッチレベルの\emph{中}でも、
どのバイナリで計算したかを特定できる。どちらも1.0.0を名乗る二つの実行の違いも、これで
追跡できる。

\section{「結果が変わりうる」とはどういうことか}
\label{sec:policy:results}

本章で最も重要な節である。リリースのたびに数値が黙って変わるソルバーは、後から結果を説明する必要がある
用途には使えない。そこで次のことを約束する。

\begin{quote}
  \textbf{同じメジャーバージョンの中では、数値出力を変える変更はすべて挙動の変更として扱う。
  変更は告知するだけでなく、測定したケースと差の大きさを示して定量的に文書化する。}
\end{quote}

この約束を守れるのは、リポジトリがすでに次の三つの場合を区別し、意図的に分けて扱っているからである。

\begin{description}

  \item[リファクタリング（ビット単位で一致すること）] 何も変えないことを意図した構造の変更である。
    求めるのは「差が小さいこと」ではなく、比較したすべてのケースのすべての列で\file{tools/compare.py}が
    \code{0.000e+00}を示すことである。これは努力目標ではなく、実際に達成できる要求である。ソルバーがそうなるように
    書かれているからである。\cref{ch:reproducibility}の決定論的な縮約、CRSの格納順序、登録順序、
    二方向の順序付けは、このために用意したものである。現在のコードに至る構造の変更も、すべてこの方法で
    検証した（\cref{sec:layout:workflow}）。

  \item[バグ修正（古い数値が誤っていたため結果が変わる）]
    この場合、差は偶然ではなく修正の目的そのものである。\cref{ch:revisions}には、何が誤っていたか、どのケースの
    結果がどれだけ変わるかを記載する。古い挙動が特定の条件でだけ誤っていた場合は、その条件も記載し、
    自分の結果が影響を受けるかどうかを判断できるようにする。この種のバグ修正はPATCHリリースで行う。

  \item[意図的な手法の変更] 離散化、既定値、収束判定などを変える場合である。結果はあらゆる箇所で
    変わり、バグ修正には当たらない。少なくともMINORの更新が必要であり、既存の入力ファイルの意味が変わるなら
    MAJORの更新が必要である。あわせて、具体的なケースで変化を定量化した項目を\cref{ch:revisions}に載せる。

\end{description}

どの場合も、確認には\file{tools/compare.py}を用いる。このツールは\file{History.dat}と最後の
\file{field*.dat}を列ごとに比較し、各列の最大差をその列のスケールで正規化する。ソルバー診断用の
四つの列（\code{NK\_iter}、\code{NK\_resid}、
\code{BCG\_iter}、\code{BCG\_resid}）は、合否の判定には含めないが報告はする。比較結果を読むときは
この点が重要である。物理量の列がすべて\code{0.000e+00}で反復回数だけが変わった場合、
変わったのは答えへのたどり着き方であり、答えそのものではない。

\begin{warn}[比較は同じ条件で行う]
  異なる線形ソルバーのバックエンドどうしや、AMGCLありとなしのビルドどうしを比べても、ビット単位の一致は
  得られず、何の証拠にもならない。バックエンドどうしは線形ソルバーの許容誤差の範囲でしか一致しない
  （\cref{sec:linsolve:choosing}）。また\val{"amgcl"}は、スレッド数を変えると再現性すらない。
  比較する両方の実行で、バックエンドとスレッド数をそろえる必要がある。
\end{warn}

\section{入力ファイルの互換性}
\label{sec:policy:compat}

\textbf{保証すること。} 1.0向けに書いた入力ファイルはどの1.0.xでも動き、すべてのキーが
本マニュアルに書いた意味を保つ。TOMLファイルだけでなく、ブロック構造の\file{.inp}ファイルも同様である。
どちらも同じ木構造に読み込まれるからである（\cref{sec:inputformat:legacy}）。1.1では省略可能なキーが新たに
加わることがあるが、書かなければ現在と同じ挙動になる。

\textbf{保証しないことと、注意が必要な理由。} 名前が変わったキーや削除されたキーは見逃されない。
ソルバーが読まない項目があると、その名前を示して実行が止まる（\cref{sec:inputformat:unknown}）。
古い綴りのキーを含む入力ファイルは、そのキーを既定値として実行されるのではなく、はっきり失敗する。
どの検査でも検出できないのは、名前はそのままで既定値や意味が変わったキーである。入力ファイルで省略した
ブロックが何を意味するかについても同じことがいえる（\cref{sec:inputformat:defaults}）。

\begin{warn}[既定値の変更は検出されないため一覧に記載する]
  パーサーはこれを警告できないので、方針で補う。\textbf{キーの名前を変えたり削除したりするリリース、
  既定値を変えるリリースでは、その内容を\cref{ch:revisions}に記載する}。新旧の綴り、または新旧の値を
  併記する。既定値の変更を見落とさない手段はこの一覧を読むことしかない。新しいリリースで改めて検証する余裕のない
  ケースを実行する前には、必ず目を通す必要がある。
\end{warn}

結果ディレクトリに含まれる入力ファイルのコピーは保管しておくこと（\cref{sec:output:dir}）。このコピーは実行に
用いた入力そのものである。後に新しいバージョンを使っていて、ある結果が実際にどの入力から計算されたかを
確かめる必要が生じたとき、拠り所になる。

\section{出力形式の安定性}
\label{sec:policy:outputs}

\file{field*.dat}は出力ファイルであると同時に、ASCII形式のリスタートファイルでもある
（\cref{sec:output:field}）。そのため、列構成は\textbf{同じメジャーバージョンの中では固定}する。
1.0.0で書いたリスタートファイルは1.0.xでも1.1でも読み込め、列の意味も変わらない。
\key{Output.Dimensioned}の設定にかかわらず有効数字18桁で書き出すのもこのためである。

\file{History.dat}と\file{CalcInfo.dat}は診断用のファイルであり、方針はもう少し緩い。\textbf{マイナーリリースで
末尾に列が追加されることがある。} どちらのファイルも、全列の名前を並べたヘッダー行で始まる。したがって次のようにする。

\begin{quote}
  \textbf{列は位置ではなく、必ずヘッダーの名前で読む必要がある。}
\end{quote}

\file{History.dat}の列を番号で参照するスクリプトは、列が追加されるまでは動くが、追加された日から黙って誤った量を
返すようになる。名前で列を探すスクリプトはこの問題を起こさない。同じメジャーバージョンの中では、列の名前が変わったり
順序が入れ替わったりすることはなく、列が追加されるだけである。

\section{保守の対象}
\label{sec:policy:support}

保守するものとしないものを簡潔に、率直に述べる。

\begin{description}
  \item[保守するもの] 三つのモデル（一次元の鉛直一様モデル、二次元の鉛直一様モデル、膜厚方向の
    濃度分布を扱う\val{"CZ"}モデル、\cref{ch:concentration}）、\cref{ch:cases}の同梱ケース、
    \cref{sec:verification:suite}のテストスイートである。これらは常にビルドし、実行し、検査する。
    このいずれかを壊す変更は、それ自体が誤った変更である。
  \item[約束しないもの] \cref{ch:notverified}に挙げた事項すべて。これらは検証が欠けていることが判明している点であり、
    検証済みと誤解されないように列挙している。今後埋めることを約束するものではない。また、未検証の量が
    誤っているわけではない。Yuzenが裏付けを提供していない量であり、利用者が自分で確かめる必要がある。
  \item[提供しないものとその理由] ILU(0)を用いる独自のBiCGStabバックエンド。開発中に性能を測定し、リリースから
    外した。1-Dでは\val{"eigen"}より四倍から六倍遅く（\cref{sec:linsolve:choosing}）、
    \file{leveling\_1d\_reg}では\val{"eigen"}の21回に対して3676回の線形反復を要した
    （\cref{sec:verification:backends}）。このバックエンド自体にも欠陥があった。一方、\val{"eigen"}は、スレッド数を
    変えてもビット単位で一致することを含め、このバックエンドが提供していたものをすべて提供する
    （\cref{ch:reproducibility,sec:verification:linalg}）。実際の計算に用いるバックエンドを選ぶときは
    \cref{sec:linsolve:choosing}を参照されたい。
\end{description}

\section{問題の報告}
\label{sec:policy:reporting}

報告に対応できるのは、ほかの人が同じ現象を再現できる場合である。再現に必要なのは次の五つであり、
どれも手元にすでにある。

\begin{enumerate}
  \item \textbf{バージョンとビルド日}。結果ディレクトリの\file{JobInfo.txt}からコピーする。「1.0」だけでは
    足りない。同じパッチレベルの二つのバイナリは、ビルド日とコンパイラで区別する。
  \item \textbf{入力ファイル}。作業ツリーにあるファイル（その後編集したかもしれない）ではなく、結果
    ディレクトリにあるコピー（\cref{sec:output:dir}）を送る。
  \item \textbf{終了メッセージ}。実行ログの末尾近くにある、実行が終了した理由を示す行である
    （\cref{tab:termination}）。ログは標準出力に出るだけでファイルには書かれないので、リダイレクトした場合に
    限り残る（\cref{sec:output:running}）。\code{Steady state.}で止まった実行と\code{Dry out.}で止まった
    実行では、報告の内容が異なる。
  \item \textbf{\file{Warnings.txt}と\file{FatalErrors.txt}に何か書かれているか}。どちらのファイルにもバナーは
    書かれないので、空になるのは何も出力されなかった場合であり、その場合に限られる。何か書かれていれば本物のメッセージである
    （\cref{sec:output:warnings}）。致命的エラーは\file{FatalErrors.txt}\emph{と}標準エラー出力の両方に
    書かれる。ファイルが残らない異常終了は、ソルバーが入力を受け付けなかった場合だけである。このときは結果ディレクトリが
    作られる前に終了するので、メッセージは標準エラー出力にしか出ない。
  \item \textbf{現象が現れる最小のケース。} 格子を粗くし、終了時刻を早め、物理項を一つずつ
    無効にして、現象が消えるところまで絞り込む。どの項を外すと現象が消えるかがわかれば、それが報告の大きな部分になる。十秒で終わるケースなら調べてもらいやすい。
\end{enumerate}

\begin{note}[スレッド数による差は常にバグである]
  \val{"eigen"}バックエンドで、スレッド数によって結果が変わるケースがあれば、それは定義上欠陥である。
  スレッド数によらず結果が同じであることは、テストスイートが確かめている主張の一つだからである
  （\cref{ch:reproducibility,sec:verification:linalg}）。二つの
  \file{compare.py}の出力を添えて報告されたい。誤りではなく想定どおりの挙動といえるのは
  \val{"amgcl"}の場合だけである。このバックエンドは内部の縮約が非決定論的だからである（\cref{sec:linsolve:choosing}）。
\end{note}
